% ID: SEL-D0066-OA-E01
% ESTADO: cotejada. Segunda pasada visual y borrador compilado revisados.
% Fuente: Selectividad/sel_2010_matematicas_aplicadas/mat_apli/mat_apli_sept/mat_apli_sept_rb/EXAMEN-ANDALUCÍA-2.pdf
% SHA-256: 24c494a5d013dbc7f5adaac9896dee2d79f22540518359c67eac29945a26d4f0
% Página: 1; opción A; ejercicio original 1.

Sean las matrices
\[
A=\begin{pmatrix}a&1\\0&2\end{pmatrix},\quad B=\begin{pmatrix}1&b\\0&3\end{pmatrix}\quad\text{y}\quad C=\begin{pmatrix}1&3\\2&5\end{pmatrix}.
\]
\begin{enumerate}[label=\alph*)]
\item \textbf{(1 punto)} Halle los valores de $a$ y $b$ para que se verifique $A-B+A\cdot B^t=C$.
\item \textbf{(0.75 puntos)} \textquestiondown{}Existe algún valor de $b$ para el que el producto $B\cdot B^t$ sea igual a la matriz nula?
\item \textbf{(0.75 puntos)} Para $a=0.5$ y $b=1$, halle la matriz $X$ que verifica la igualdad $A\cdot X+B=O$, ($O$ representa la matriz nula).
\end{enumerate}
